% ============================================================
% Appendix A --- Team AI Transparency Log (rubric dimension E)
% Included from Method.tex after \appendix.
% TODO (team): fill in the exact product versions/dates in Table A.1.
% ============================================================
\section{AI Transparency Log}
\label{app:ai_log}

This appendix records the distance between raw generative-AI output and the
claims that survived into this report; it is a record of auditing, not a count
of prompts. AI tools were used throughout the project in the permitted way.
Every claim listed below was re-derived, re-run against the team's own code, or
falsified before it was allowed into the text, and three of the audits changed
a claim that an earlier draft had already stated. Conversation transcripts are
retained by the team and can be supplied on request; they are not appended
here.

\subsection*{A1. AI tools used}

\begin{table}[H]
\centering\small
\caption{Generative-AI tools used in Project 2. Entries describe the use that
actually reached the report; routine autocomplete and spelling checks are
covered by the last row.}
\label{tab:ai_tools}
\begin{tabular}{@{}p{3.9cm}p{3.3cm}p{7.4cm}@{}}
\toprule
Tool & Version / date & Primary use \\
\midrule
DeepSeek-based assistant in DeepSeek Harness (\texttt{deepseek-flash})
  & current build, Oct--Nov 2026
  & Drafting and re-deriving \S1.2 and \S2.1--\S2.4 (model, weak form,
    truncation error, interface conductance, boundary elimination, discrete
    structure); writing the independent verification scripts in
    \path{verify/} and \path{test/}; adversarial review of each other's text. \\
Claude (Anthropic), web application
  & current build, Oct--Nov 2026
  & Drafting of the algorithm section, integration notes for merging the
    sections, figure and table captions, wording and consistency passes. \\
ChatGPT (OpenAI) and Gemini (Google)
  & current builds, Oct--Nov 2026
  & \LaTeX{} error diagnosis, document structure, phrasing and completeness
    checks; carried over from the tool set recorded in the Project 1 log. \\
\midrule
All of the above, routine use
  & ---
  & Autocomplete, spelling, reference formatting, table alignment. No numerical
    value was accepted without reproduction (see A3). \\
\bottomrule
\end{tabular}
\end{table}

\subsection*{A2. Significant AI-assisted claim audits}

Three interactions are recorded below because each one changed a mathematical
or numerical claim rather than simply rewording it. The common failure pattern
is a claim that is true for the face coefficient, or for one dimension, being
silently promoted to the solution of the two-dimensional problem.

\begin{table}[H]
\centering\small
\caption{Claim audits. Column 3 states what would have shown the claim to be
wrong.}
\label{tab:ai_audits}
\begin{tabular}{@{}p{0.95cm}p{7.4cm}p{5.8cm}@{}}
\toprule
Item & Prompt / claim / assumptions & Risk and possible falsifier \\
\midrule
1 &
``Arithmetic averaging of the interface conductance produces an
$\mathcal{O}(1)$ error in the interface flux, so the arithmetic-averaged
solution is $\mathcal{O}(1)$ inaccurate.''
\textbf{Assumption hidden in the wording:} the face-coefficient error and the
error of the discrete solution are the same object. &
The claim mixes three different quantities: the relative face-conductance
error $(1-c)^2/(4c)$ (genuinely $\mathcal{O}(1)$ in $h$), the nodal error and
the midline interface-flux error (first order in $h$), and the \emph{maximum}
interface-flux error over the square, which is $\mathcal{O}(1)$ for a different
reason --- the exact Dirichlet trace on $y=0,1$. \textbf{Falsifier:} a
refinement study. If the nodal error did not halve under mesh refinement, the
first-order part of the claim would be wrong. \\
2 &
``Since $A_h \succeq a_{\min}A_h^{(1)}$ in the Loewner order and both matrices
are $M$-matrices, the inverses compare entrywise,
$A_h^{-1}\le a_{\min}^{-1}(A_h^{(1)})^{-1}$; hence
$\|A_h^{-1}\|_\infty\le 1/(8a_{\min})$.'' \textbf{Assumption hidden:}
spectral order implies entrywise order. &
The implication is false: monotonicity of $M$-matrix inverses needs an
\emph{entrywise} hypothesis, while growing $a$ makes the off-diagonals
$-a_f/h^2$ \emph{more negative}. \textbf{Falsifier:} evaluate both inverses for
a strongly varying coefficient, e.g.\ $a=1+9\sin^2(\pi x)$ at $n=63$, where the
entrywise ratio reaches $3.10>1$. The conclusion
$\|A_h^{-1}\|_\infty\le1/(8a_{\min})$ is nevertheless true and must be proved by
a barrier argument instead. \\
3 &
``The arithmetic operator is symmetric and its shared face flux is
conservatively balanced, so the interface physics is correct.'' &
Symmetry and local conservation are necessary but not sufficient; they are
satisfied by a matrix whose interface flux is wrong by $36\%$.
\textbf{Falsifier:} compare the computed interface flux with the exact
$J_*=-20/11$ on a benchmark whose exact solution is known. It differs by
$0.662$ while the two-sided reconstructed jump is $3.85\times10^{-13}$ and the
symmetry defect is zero. \\
\bottomrule
\end{tabular}
\end{table}

\subsection*{A3. Independent verification and decision}

\begin{table}[H]
\centering\small
\caption{What was checked, where the evidence lives, and what changed.}
\label{tab:ai_verification}
\begin{tabular}{@{}p{0.95cm}p{6.2cm}p{7.0cm}@{}}
\toprule
Item & Independent check and evidence location & Final decision / modification / limitation \\
\midrule
1 &
Re-derived the truncated-series argument in one dimension, where the discrete
system reduces to a tridiagonal chain with a closed-form flux; reproduced the
leading term $h(a_L-a_R)^2/(a_L+a_R)^2$ analytically and confirmed it against
$N=31,\dots,255$ (\path{verify/check_interface_order.py}, Table~\ref{tab:interface_refinement}).
Then re-assembled the full two-dimensional layered problem and compared the
nodal, midline and maximal face-flux errors with \S4.3. &
The two-layer statement of \S2.2 and Remark~\ref{rem:arith_2d_boundary} replaced
the single blanket $\mathcal{O}(1)$ claim: first-order nodal and midline flux
error, $\mathcal{O}(1)$ maximal face-flux error near the prescribed horizontal
boundaries, $\mathcal{O}(1)$ face-conductance error. \textbf{Limitation:} the
maximum is a boundary effect of the Dirichlet square; we state it as an
observation of this benchmark, not as an estimate valid for other boundary
data. \\
2 &
Ran the counterexample search and the bound check in
\path{verify/check_stability.py}: entrywise ratio for $a=1+9\sin^2(\pi x)$,
$\|A_h^{-1}\|_\infty$ against $1/(8a_{\min})$ for five coefficients in one and
two dimensions, and the comparison with the constant-coefficient maximum
$0.0736<\frac18$ (\path{verify/check_stability.py} output,
Remark~\ref{rem:entrywise_false}, Remark~\ref{rem:conv_constant}). &
Proof of Theorem~\ref{thm:convergence} rewritten with a maximum-principle
barrier: $A_h\bm\varphi=\bm 1$ exactly for
$\varphi=\frac12x(1-x)/a_{\min}$ extended by its own boundary values, and
$A_h(\bm\varphi-\bm w)=0$ with $\bm\varphi-\bm w\ge0$ on $\partial\Omega$.
\textbf{Limitation:} equality in $\|\bm w\|_\infty\le1/(8a_{\min})$ holds only
in one dimension; in two the constant is an upper bound ($0.0736$ versus
$0.125$). \\
3 &
Separated the three diagnostics on the same computed solution: symmetry defect
$\|A_h-A_h^T\|/\|A_h\|$, reconstructed two-sided flux jump, and error against
the exact flux (\S\ref{sec:results} and the interface diagnostics of
\S\ref{sec:algorithm_impl}; independent
re-computation in \path{test/validate_interface.py}). &
Remark~\ref{rem:symmetry_vs_accuracy} and \S4.3 now report the three numbers
side by side instead of treating symmetry as evidence of physical accuracy.
\textbf{Limitation:} the jump diagnostic is only defined where a two-sided
reconstruction exists, so it is not a substitute for an exact solution. \\
\bottomrule
\end{tabular}
\end{table}

Two further routine checks are recorded for completeness. The smooth benchmark
is \emph{blind} to the $\theta$-correction of the face coefficient when
$a=1+x$: arithmetic and midpoint averaging then produce entrywise identical
matrices, so the second-order table in \S4.2 cannot validate the correction
term (Remark~\ref{rem:theta_blind}). And no convergence claim is made anywhere
in this report without a named norm and a four-level refinement table; claims
about the layered benchmark are exactness statements
(Theorem~\ref{thm:nodal_exactness}), not order statements, because
$u_*\notin H^2(\Omega)$ there (Remark~\ref{rem:thm2_scope}).

\subsection*{A4. Intellectual ownership declaration}

The mathematical formulation, the discretisation, the solvers, the experiments,
the validation evidence and the conclusions are the team's own work.
Generative AI was used for drafting, clarification and preliminary checking, and
every AI-assisted claim was audited as recorded above. Each member affirms that
they can explain and modify the team-owned discretisation and solver code,
justify the mathematical claims in their own sections, and state the
assumptions and the tests applied at each library interface, without AI
assistance during the viva. The viva may ask for a small live modification or
diagnostic on team-owned work.

\begin{table}[H]
\centering\small
\caption{Ownership declaration. Signatures are collected on the submitted PDF.}
\label{tab:ai_ownership}
\begin{tabular}{@{}p{2.3cm}p{2.2cm}p{4.6cm}p{4.9cm}@{}}
\toprule
Name & Signature & Team-owned artefact / contribution & Trusted dependency and interface \\
\midrule
Tianrui Ji & \underline{\hspace{1.6cm}} & \path{src/discretisation.py}, \path{src/solvers.py}, \path{run_all.py}, \S3, Figures 1--6 & \texttt{scipy.sparse} CSR assembly and \texttt{scipy.sparse.linalg} CG/LU: checked by solving the same system twice, CG against SuperLU, residual $\le10^{-12}$ \\
Meichen Ding & \underline{\hspace{1.6cm}} & \S1.2, \S2.1--\S2.4, Remarks on the $\theta$-correction, mesh parity, two-dimensional boundary and symmetry, Appendix~A & \texttt{numpy.linalg} dense solves and \texttt{scipy.sparse.linalg.eigsh} used only for corroboration; checked against explicit residuals and against the exact layered solution \\
Jingyan Wei & \underline{\hspace{1.6cm}} & \S2.4 variational $\mathcal P_1$ framework and matched-settings protocol, interface extension (\S3.3), Appendix~A audit 2 & \texttt{scipy.sparse} sparse factors and eigensolvers at the same tolerance as the FDM comparison; interface conventions tested on the manufactured solution \\
Han Liu & \underline{\hspace{1.6cm}} & \S5, Appendix~C, integration of the section files, tables, figures and references into one report; repository packaging & \LaTeX{} and \texttt{latexmk} build: clean build reproduced from a fresh clone (zero errors, zero undefined references) \\
Jiasheng Xu & \underline{\hspace{1.6cm}} & \path{test/test_validation.py}, \path{test/validate_smooth.py}, \path{test/validate_interface.py}, \path{test/compare_upstream.py}, \S4.1 & \texttt{unittest} and \texttt{numpy} reference solves independent of the production assemblers; the strict-mode run is recorded with its non-zero exit rather than suppressed \\
\bottomrule
\end{tabular}
\end{table}
